\subsubsection*{Lines in three-dimensional space}

The same point-plus-direction construction extends to $\mathbb R^3$ without
changing its meaning. This continuity matters: a good definition should survive
a change of dimension. At the same time, three-dimensional space creates a new
possibility. Two nonparallel lines need not intersect; they may be skew.

The section teaches the student to separate questions that coincide in the plane.
Parallel direction, perpendicular direction, intersection, and perpendicular
lines are different conditions in space. A zero dot product between direction
vectors gives a right angle between directions, but the geometric lines are
perpendicular only if they also meet. A zero cross product detects parallel
directions, but additional information is needed to decide whether the lines
coincide or are distinct.

The student should be able to construct a spatial line from a point and direction
or from two points, use parametric and symmetric descriptions, solve for
intersections, distinguish parallel, intersecting, perpendicular, and skew lines,
and explain why three scalar equations for two parameters may be inconsistent.
The deeper lesson is that increasing dimension changes which geometric outcomes
are possible even when the algebraic language remains similar.
